Eine zentrale Frage in der Analysis ist es zu verstehen, wann eine Gleichung der Form \( f(x,y) = 0 \) nach einer Variablen (z.B. \( y \)) aufgelöst werden kann, sodass eine Funktion \( g(x) \) existiert mit 


\[
f(x, g(x)) = 0.
\]


Die Menge der Lösungen wird dann als Graph der Funktion \( g \) dargestellt.

\begin{beispielbox}
\begin{enumerate}
    \item Sei \( f(x,y) = y^2 - x^2 \). Dann folgt aus der Gleichung \( f(x,y) = 0 \), dass \( g(x) = \pm x \). 
    \item Betrachten wir \( f(x,y) = y^2 - x^4 \). Nach \( y \) auflösen liefert \( g(x) = \pm x^2 \), aber eine Auflösung nach \( x \) ist global möglich mit \( x = \pm \sqrt{y} \).
    \item Sei \( f(x,y) = x^2 + y^2 - 1 \). Diese Gleichung beschreibt einen Kreis. Für \( x = 0 \) ist eine Auflösung nach \( y \) möglich: \( y = \pm 1 \), für \( y=0 \) folgt \( x = \pm 1 \).
\end{enumerate}
\end{beispielbox}

\begin{bemerkungbox}
Diese Beispiele zeigen, dass eine solche Umformung nicht immer für beliebige Punkte möglich ist. Die zentrale Idee besteht darin, das **Theorem der impliziten Funktion** zu beweisen.
\end{bemerkungbox}